\documentclass[preprint,12pt]{elsarticle}

\usepackage{amsmath,amssymb,amsfonts,amsthm}
\usepackage{graphicx}
\usepackage{hyperref}
\usepackage{booktabs}

\newtheorem{theorem}{Theorem}

\newcommand{\SU}{\mathrm{SU}}
\newcommand{\UU}{\mathrm{U}}
\newcommand{\Nf}{N_{\!f}}
\newcommand{\Nc}{N_{\!c}}
\newcommand{\Tr}{\mathrm{Tr}}
\newcommand{\adj}{\mathrm{adj}}
\newcommand{\diag}{\mathrm{diag}}
\newcommand{\MeV}{\,\mathrm{MeV}}
\newcommand{\GeV}{\,\mathrm{GeV}}
\newcommand{\TeV}{\,\mathrm{TeV}}

\begin{document}

\begin{frontmatter}

\title{The quark mass waterfall from $\mathcal{N}{=}2$ interleaving}

\author{XXXXX}
\address{XXXXX}

\begin{abstract}
  We show that the empirical ``waterfall'' pattern in quark masses---the
  fact that consecutive triples $(t,b,c)$, $(b,c,s)$, $(c,s,0)$,
  $(s,0,\chi)$ (where $\chi \equiv m_u + m_d$) each approximately satisfy the Koide ratio $K = 2/3$
  with appropriate sign assignments---is a structural
  consequence of the doublet structure that emerges from
  confinement of the $\SU(3)_F$ family gauge theory.
  The confined spectrum naturally produces two mass eigenvalues
  per family---a direct branch $a_i = \mu q_i^2$ (from the
  adjugate $\adj(M)_i$) and an
  inverse branch $d_i = \nu/q_i^2$ (from the meson $M_i$)---whose
  pairing is electroweak isospin.
  The two orderings are opposite: the direct masses
  increase with charge, but the inverse masses \emph{decrease}.
  We prove an \emph{interleaving theorem}: the six unmixed masses
  automatically form a descending ladder
  $d_1 > a_3 > d_2 > a_2 > d_3 > a_1$, alternating between branches,
  whenever the Seiberg scale $\nu/\mu$ lies in a calculable window.
  The Standard Model quark masses satisfy this window condition.
  The sign flip in the diagnostic $K(b,c,-s)$ signals
  the branch mismatch within a mixed window, though its
  detailed assignment requires the quartic mixing.
  The interleaving requires no new charges beyond the three family
  eigenvalues; the waterfall is the observable signature of
  the composite doublet structure producing the up-type quark sector
  alongside the down-type sector from a single gauge theory.
  The identification of specific SM quarks with specific branch
  entries relies on the companion letter's branch assignments.
\end{abstract}

\end{frontmatter}

%----------------------------------------------------------------------
\section{Introduction}

The charged-lepton masses satisfy the Koide ratio
$K \equiv \sum m_i / (\sum\sqrt{m_i})^2 = 2/3$ to $0.001\%$,
and the inverse down-quark masses satisfy
$K^{-1}(d,s,b) = 2/3$ to $0.3\%$~\cite{Koide:1983qe,Rivero:2005}.
In the companion letter~\cite{letter_q2}, we showed that both
relations follow from a quadratic charge--mass law
$m_i = \mu\,q_i^2$ and its Seiberg dual $\hat{m}_i \propto 1/q_i^2$
in a confined $\SU(3)_F$ gauge theory with three flavours.

To summarise the companion letter's framework briefly:
the family symmetry $\SU(3)_F$ confines at a scale~$\Lambda$,
producing mesons $M_{ij} = Q^i\tilde{Q}_j$ and their adjugates
$\adj(M)_{ij}$.  Each family's charge $q_i = z_0 + z_i$
enters a quadratic mass law $m_i = \mu\,q_i^2$ (direct branch)
or $\hat{m}_i = \nu/q_i^2$ (inverse branch).
The Koide identity $K = 2/3$ is a trace identity
of the charge operator~$H_Q = Q_F^\dagger Q_F$, valid for
the direct branch at any direction in the Cartan plane.

There is a further empirical observation.
If one lists all six quark masses in descending order
$(m_t, m_b, m_c, m_s, m_u, m_d)$ and evaluates the Koide ratio
on sliding windows of three consecutive masses, the results cluster
near $K = 2/3$ with suitable sign assignments.
This ``waterfall'' pattern was noted in
Refs.~\cite{Rivero:2005,Zenczykowski:2012}
and involves mixed quark flavours (up-type and down-type together),
unlike the pure-sector relations.
Here we show that the waterfall is a structural consequence
of the composite doublet structure of the confined family
gauge theory:
the direct and inverse mass branches interleave automatically,
producing a descending ladder whose consecutive windows
inherit approximate $K = 2/3$ from the underlying charge
normalisation.

\paragraph{Conventions.}
We work throughout in the charge-squared language:
masses are $m_i = \mu\,q_i^2$ (direct) or $\hat{m}_i = \nu/q_i^2$
(inverse), with $q_i = z_0 + z_i$ the family charges
($z_0 = 1/\sqrt{6}$, $\sum z_i = 0$, $\sum z_i^2 = 1/2$).
We write $\chi \equiv m_u + m_d$ for the lightest isospin sum.
Square roots $\sqrt{m_i}$ appear only as derived diagnostics
inside the ratio~$K$, never as primitive variables.

%----------------------------------------------------------------------
\section{The two branches}
\label{sec:branches}

The family charge operator is
$Q_F = z_0 I + T$, with $T$ in the Cartan of $\SU(3)_F$.
Its square $H_Q = Q_F^\dagger Q_F = \diag(q_1^2, q_2^2, q_3^2)$
has trace $\Tr(H_Q) = \sum q_i^2 = 1$ by the normalisation
conditions.  The mass ansatz $m_i = \mu\,q_i^2$ gives
\begin{equation}
  K \;=\; \frac{\sum q_i^2}{(\sum q_i)^2}
  \;=\; \frac{1}{3/2} \;=\; \frac{2}{3}\,,
  \label{eq:K_trace}
\end{equation}
a trace identity of $H_Q$, independent of the charge direction~$\alpha$.
In confined $\SU(3)$ SQCD with $\Nf = \Nc = 3$~\cite{Seiberg:1994,CSS:1996,IS:1995},
the quantum-modified constraint $\det M = \Lambda^6$ and
$F$-flatness give the meson VEVs $\langle M_i \rangle \propto 1/m_i
\propto 1/q_i^2$.  Defining the inverse-branch scale
$\nu$ (whose relation to $\Lambda$ and $\mu$ is derived
in the companion letter; for the present analysis
only the parametric form matters),
the inverse-branch masses are
\begin{equation}
  \hat{m}_i \;=\; \frac{\nu}{q_i^2}\,.
  \label{eq:inverse}
\end{equation}
Since $a_i\,d_i = \mu q_i^2 \cdot \nu/q_i^2 = \mu\nu$,
the geometric mean $\sqrt{a_i\,d_i} = \sqrt{\mu\nu}$ is
\emph{family-independent}---a direct consequence of
$\det M = \Lambda^6$.

Unlike the direct branch, $K(\hat{m}_1, \hat{m}_2, \hat{m}_3)$
is \emph{not} automatically $2/3$: the reciprocal charges
$1/q_i$ do not satisfy the same normalisation as the~$q_i$.
The empirical inverse-Koide relation
$K^{-1}(d,s,b) \approx 2/3$ (to $0.3\%$)~\cite{Rivero:2005}
is a non-trivial constraint---it implies that the down-quark
sector has its own effective charges satisfying
the trace condition, related to the lepton charges by
a small angle shift~$\Delta\alpha$~\cite{letter_q2}.
(The notation ``$K^{-1}$'' denotes $K$ evaluated on
reciprocal masses $1/m_i$.)

The key structural point is the \emph{ordering}.
Since $q_1 < q_2 < q_3$ (convention):
\begin{alignat}{3}
  &\text{Direct:}\quad  & a_1 < a_2 < a_3 \qquad
  & (a_i = \mu\,q_i^2, \text{ increases with } q)\,,
  \label{eq:direct_order} \\
  &\text{Inverse:}\quad & d_1 > d_2 > d_3 \qquad
  & (d_i = \nu/q_i^2, \text{ decreases with } q)\,.
  \label{eq:inverse_order}
\end{alignat}
The two orderings are \emph{opposite}.
This is the algebraic origin of the waterfall.

%----------------------------------------------------------------------
\section{The emergent doublet structure}
\label{sec:N2}

The $\mathcal{N}{=}1$ confined theory already contains both
branches---direct ($\adj(M)_i \propto q_i^2$) and inverse
($M_i \propto 1/q_i^2$)---but treats them independently.
The composite spectrum, however, has a natural doublet structure:
for each family~$i$, the pair $(M_i, \adj(M)_i)$ furnishes
two mass eigenvalues with opposite charge dependence.
This pairing \emph{is} electroweak isospin---the down-type
and up-type composites within each generation---and has the
mathematical form of an $\mathcal{N}{=}2$ hypermultiplet
$(Q_i, \tilde{Q}_i)$ under the $\SU(2)_R$ that rotates them.
The ``$\mathcal{N}{=}2$ structure'' is thus emergent from
confinement plus isospin, not a UV completion that is subsequently
broken~\cite{APS:1997,SeibergWitten:1994a}.

\paragraph{The effective quartic.}
The coupling between the two branches can be parametrised by a
quartic meson interaction.
Formally, this quartic is identical to what one obtains by
introducing an auxiliary adjoint chiral field~$\Phi$ with
mass~$\mu_\Phi$ and the standard $\mathcal{N}{=}2$
coupling $W \supset \sqrt{2}\,g\,\tilde{Q}\,\Phi\,Q$,
then integrating $\Phi$ out via its $F$-term~\cite{APS:1997}.
The result is an effective $\mathcal{N}{=}1$
superpotential:
\begin{equation}
  W_{\rm eff} \;=\;
  -\frac{g^2}{\mu_\Phi}\Bigl[\Tr(Q\tilde{Q})^2
  - \tfrac{1}{3}\bigl(\Tr Q\tilde{Q}\bigr)^2\Bigr]
  \;+\; m_i\,Q^i\tilde{Q}_i\,.
  \label{eq:Weff}
\end{equation}
We emphasise that $\Phi$ here is a bookkeeping device for
the quartic, not a physical UV field: the confined theory
is defined at the family confinement scale
$\Lambda_F \sim \mathcal{O}(100)\MeV$, and the
``$\mathcal{N}{=}2$ breaking scale'' $\mu_\Phi$ parametrises
the strength of the inter-branch coupling in the
composite dynamics.

\paragraph{The decoupling limit.}
As $\mu_\Phi \to \infty$, the quartic vanishes, the branches
fully decouple, and one recovers pure $\mathcal{N}{=}1$
SQCD with masses~$m_i$.
For finite~$\mu_\Phi$, the quartic leaves a calculable
residual coupling between the direct and inverse branches.

\paragraph{The residual coupling in the confined phase.}
In the confined phase, $Q\tilde{Q} \to M$ (the meson).
The quartic becomes
\begin{equation}
  W_{\rm quartic} \;=\;
  -\frac{g^2}{\mu_\Phi}\Bigl[\Tr M^2
  - \tfrac{1}{3}(\Tr M)^2\Bigr],
  \label{eq:quartic}
\end{equation}
which modifies the $F$-term equations by terms
of order $g^2 M_i/\mu_\Phi$.  Since $M_i \propto 1/q_i^2$,
the fractional correction to the inverse branch is
\begin{equation}
  \frac{\delta M_i}{M_i}
  \;\sim\; \frac{g^2\,M_i}{\mu_\Phi}
  \;\propto\; \frac{1}{m_i}\,,
  \label{eq:correction}
\end{equation}
which is \emph{smallest} for the heaviest family
(largest $m_i$) and \emph{largest} for the lightest.
The waterfall windows involving heavier quarks are
therefore expected to be closer to $K = 2/3$ than
those dominated by lighter quarks, consistent with the
pattern in Table~\ref{tab:waterfall}.

\paragraph{The up/down pairing.}
In the confined theory, each family~$i$ produces two
composite operators: $M_i$ (meson) and $\adj(M)_i$
(adjugate).  These give two independent mass parameters:
$a_i = \mu\,q_i^2$ (direct, from the adjugate) and
$d_i = \nu/q_i^2$ (inverse, from the meson VEV).
The pairing of a direct mass with an inverse mass
within each family---i.e.\ the pairing of a
down-type quark with an up-type quark---is the
composite manifestation of electroweak isospin.%
\footnote{The involution
$M_i \leftrightarrow \adj(M)_i = \Lambda^6/M_i$,
which exchanges $q_i^2 \leftrightarrow 1/q_i^2$
within each generation, is a Seiberg-duality
transformation between the electric and magnetic
descriptions of the confined phase;
see the companion paper~\cite{letter_q2} for details.}

\paragraph{Historical note.}
That the MSSM requires two Higgs doublets was
first shown by Fayet~\cite{Fayet:1976,Fayet:1977};
the observation that $H_u$ and $H_d$ can be identified
with the two halves of a composite doublet is a natural
consequence of confinement with $\Nf = \Nc$.
What is new here is that the \emph{mass spectrum}---not
only the Higgs sector---carries the doublet imprint:
the direct and inverse branches, their opposite orderings,
and their interleaving into the waterfall are all consequences
of having two composite operators per family.
The connection to electroweak symmetry via Seiberg
duality of $\SU(4)$ SQCD~\cite{CST:2011} is discussed
in the companion letter.

%----------------------------------------------------------------------
\section{The interleaving theorem}
\label{sec:interleaving}

In the $\varepsilon \to 0$ limit, the six masses are
$\{a_1, a_2, a_3, d_1, d_2, d_3\}$ with the orderings
(\ref{eq:direct_order})--(\ref{eq:inverse_order}).

\begin{theorem}[Interleaving]
The six unmixed masses interleave as
\begin{equation}
  d_1 > a_3 > d_2 > a_2 > d_3 > a_1
  \label{eq:ladder}
\end{equation}
if and only if
\begin{equation}
  \max\!\bigl(q_1^2 q_3^2,\; q_2^4\bigr)
  \;<\; \frac{\nu}{\mu}
  \;<\; q_2^2\,q_3^2\,.
  \label{eq:window}
\end{equation}
\end{theorem}

\begin{proof}
The five adjacent inequalities reduce to three independent conditions:
\begin{alignat}{2}
  d_1 > a_3 &\iff \nu/q_1^2 > \mu q_3^2
  &&\iff \nu/\mu > q_1^2 q_3^2\,,\\
  a_3 > d_2 &\iff \mu q_3^2 > \nu/q_2^2
  &&\iff \nu/\mu < q_2^2 q_3^2\,,\\
  d_2 > a_2 &\iff \nu/q_2^2 > \mu q_2^2
  &&\iff \nu/\mu > q_2^4\,.
\end{alignat}
The remaining two ($a_2 > d_3$, $d_3 > a_1$) are implied by
these (since $q_1 < q_2 < q_3$).
The combined condition is~(\ref{eq:window}).
Since $q_1^2 q_3^2 < q_2^2 q_3^2$ and $q_2^4 < q_2^2 q_3^2$
(both using $q_2 < q_3$), the window is non-empty.
\end{proof}

The ladder~(\ref{eq:ladder}) alternates strictly between
inverse and direct branches:
\begin{equation}
  \underbrace{d_1}_{\text{inv}},\;
  \underbrace{a_3}_{\text{dir}},\;
  \underbrace{d_2}_{\text{inv}},\;
  \underbrace{a_2}_{\text{dir}},\;
  \underbrace{d_3}_{\text{inv}},\;
  \underbrace{a_1}_{\text{dir}}\,.
  \label{eq:alternating}
\end{equation}
The seesaw structure is manifest: the \emph{heaviest}
mass ($d_1$) comes from the \emph{inverse} branch acting on the
\emph{lightest} charge ($q_1$), while the \emph{lightest}
mass ($a_1$) comes from the \emph{direct} branch acting on
the same charge.

%----------------------------------------------------------------------
\section{Numerical verification}
\label{sec:numerics}

We evaluate the interleaving condition using the
physical charge values.
The lepton angle $\alpha_{\rm lep} = 0.2222$~rad
gives~\cite{letter_q2}
\begin{equation}
  q_1 = 0.01648\,,\quad
  q_2 = 0.23686\,,\quad
  q_3 = 0.97140\,.
  \label{eq:charges}
\end{equation}
Following the companion letter's branch assignments
(top on inverse, bottom on direct),
we fix $\mu = m_b/q_3^2 = 4430\MeV$ and
$\nu = m_t q_1^2 = 46.8\MeV$.  The window
condition~(\ref{eq:window}) reads
\begin{equation}
  \max(2.56 \times 10^{-4},\; 3.15 \times 10^{-3})
  < \underbrace{\frac{\nu}{\mu}}_{= 1.06 \times 10^{-2}}
  < 5.29 \times 10^{-2}\,,
\end{equation}
which is satisfied.  The resulting ladder is:

\begin{table}[h]
\centering
\caption{The interleaved mass ladder from 3~families $\times$ 2~branches.
  All masses in MeV.  SM identifications are tentative for the
  intermediate entries.}
\label{tab:ladder}
{\small
\begin{tabular}{llrl}
  \toprule
  Label & Branch & Mass [MeV] & SM candidate \\
  \midrule
  $d_1 = \nu/q_1^2$ & inverse & $172\,400$ & top \\
  $a_3 = \mu q_3^2$ & direct  & $4\,180$   & bottom \\
  $d_2 = \nu/q_2^2$ & inverse & $835$      & charm (shifted) \\
  $a_2 = \mu q_2^2$ & direct  & $249$      & strange (shifted) \\
  $d_3 = \nu/q_3^2$ & inverse & $50$       & $\chi$ (tail) \\
  $a_1 = \mu q_1^2$ & direct  & $1.2$      & $\approx 0$ \\
  \bottomrule
\end{tabular}}
\end{table}

The top and bottom entries match the SM by construction
($\mu$ and $\nu$ are fixed from them).
The intermediate entries $d_2 = 835\MeV$ and $a_2 = 249\MeV$
should be compared with $m_c = 1273\MeV$ and $m_s = 93.5\MeV$.
The sum rule provides a non-trivial test:
\begin{equation}
  T_2 = \mu q_2^2 + \frac{\nu}{q_2^2}
  = 249 + 835 = 1084\MeV\,,
  \label{eq:sum_rule}
\end{equation}
to be compared with $m_c + m_s = 1366\MeV$ (a $21\%$ tension).
The individual entries fare worse:
$d_2/m_c = 0.66$ and $a_2/m_s = 2.66$.
The residual quartic~(\ref{eq:quartic}) can redistribute
$T_2$ between the two eigenvalues, but these $\mathcal{O}(1)$
individual shifts indicate that quantitative agreement
at the per-quark level requires either
the quark-sector angle shift (option~ii of the Discussion)
or a non-perturbative treatment of the quartic mixing.
The ordering $d_2 > a_2$ (charm above strange) is robust
since it depends only on the interleaving
condition~(\ref{eq:window}), not on the individual values.

The waterfall diagnostic on the physical quark masses is:

\begin{table}[h]
\centering
\caption{The waterfall: Koide ratio for sliding windows
  of three consecutive quark masses.  Signs refer to the
  argument of $K$:
  $K(\pm\sqrt{x_1}, \pm\sqrt{x_2}, \pm\sqrt{x_3})$.
  PDG~2024 values~\cite{PDG:2024}.}
\label{tab:waterfall}
{\small
\begin{tabular}{llcc}
  \toprule
  Window & Signs & $K$ & $\delta K$ \\
  \midrule
  $(t, b, c)$ & $(+,+,+)$ & $0.6692$ & $+0.37\%$ \\
  $(b, c, -s)$ & $(+,+,-)$ & $0.6747$ & $+1.20\%$ \\
  $(c, s, 0)$ & $(+,+,+)$ & $0.6646$ & $-0.31\%$ \\
  $(s, 0, \chi)$ & $(+,+,+)$ & $0.6649$ & $-0.27\%$ \\
  \bottomrule
\end{tabular}}
\end{table}

\noindent The last two windows use the zero-mass limit:
$K(a, b, 0) = 2/3$ if and only if
$a/b = (2 \pm \sqrt{3})^2 = R^{\pm 1}$,
where $R = (2+\sqrt{3})^2 \approx 13.928$.
The ratios $m_c/m_s = 13.63$ ($R$: $-2.1\%$) and
$m_s/(m_u{+}m_d) = 13.68$ ($R$: $-1.8\%$) are both
close to~$R$.

%----------------------------------------------------------------------
\section{The sign flip}
\label{sec:sign}

The second window in the
ladder~(\ref{eq:alternating}) is
$(a_3, d_2, a_2)$---two direct entries and one inverse.
In the Koide diagnostic, $\sqrt{m}$ maps to:
\begin{equation}
  \sqrt{a_j} = \sqrt{\mu}\,q_j\,,\qquad
  \sqrt{d_j} = \frac{\sqrt{\nu}}{q_j}\,.
  \label{eq:sqrt_branches}
\end{equation}
In a pure-branch window, all three $\sqrt{m_k}$ scale the
same way with the charges and the sum is controlled by
the normalisation $\sum q_k = \sqrt{3/2}$.
In a mixed window the two branches contribute different
functions of the charges; a relative sign is needed
to obtain $K \approx 2/3$.

\paragraph{The sign assignment is empirical.}
The observed sign pattern---negative on $s$ in
$K(b, c, -s)$---does not coincide with the
minority-branch member ($d_2 = $ charm, inverse).
A simple ``mark the inverse member'' rule would
predict $K(b, -c, s)$, which gives $K \approx 3.7$.
The correct sign assignment therefore carries
information beyond the zeroth-order branch
identification; it likely reflects the
quartic redistribution~(\ref{eq:quartic}) that
shifts the unmixed values toward the physical
masses, reassigning the effective branch character
of the intermediate entries.
We regard the sign pattern as a constraint on the
mixing, not a prediction of the present analysis.

%----------------------------------------------------------------------
\section{Discussion}
\label{sec:discussion}

\paragraph{What the interleaving achieves.}
The theorem establishes that 3~families $\times$ 2~branches
produce a 6-rung descending ladder, with the seesaw
inversion $d_1 = \nu/q_1^2 \gg a_3 = \mu q_3^2$
identifying the top quark as the \emph{inverse} of the lightest
family charge.  The interleaving extends the pure-sector
Koide relations to the full quark spectrum by mixing
both branches in each window.  The $\mathcal{N}{=}1$ theory already contains both branches
(direct for leptons, inverse for down quarks in the companion
letter); the emergent doublet structure adds the
\emph{residual quartic coupling}~(\ref{eq:quartic}) between
the branches, which controls the mixing and deviations
in the waterfall windows.

\paragraph{The sum-rule tension.}
The $21\%$ mismatch of $T_2$ (Eq.~\ref{eq:sum_rule})
is not a failure of the interleaving---which is
about ordering, not absolute values---but a measure of
the residual quartic~(\ref{eq:quartic}).
Two resolutions are possible:
\begin{enumerate}
\item The quartic redistribution at finite $g^2/\mu_\Phi$
  shifts the eigenvalues away from the unmixed $a_i$, $d_i$.
  The correction~(\ref{eq:correction}) is family-dependent
  through $M_i \propto 1/q_i^2$.
\item A quark-sector angle $\alpha_q \neq \alpha_{\rm lep}$
  (1~extra parameter), changing the charge ratios.
  This option has independent support from the CKM
  angle shift $\Delta\alpha = \arctan(\pi/84)$~\cite{letter_q2}.
\end{enumerate}
Option~(ii) is more economical and testable:
it predicts that the quark waterfall and the CKM matrix
are controlled by the \emph{same} angle offset.

\paragraph{Why the constraint surface alone is insufficient.}
The Koide condition $K = 2/3$ for a mixed window
$(d_i, a_j, d_k)$ or $(a_j, d_k, a_\ell)$ reduces to a
quadratic equation in $r = \sqrt{\nu/\mu}$, whose coefficients
depend only on the charges.  There are four windows but only
one free parameter~$r$, so the system is generically overconstrained.
Explicit computation confirms that no value of~$r$ satisfies
$K = 2/3$ in all four windows simultaneously
(the minimum of $\max_w |\delta K_w|$ is~$\sim\!7\%$).
In particular, the pure-sector identity $\sum q_i^2 / (\sum q_i)^2 = 2/3$
guarantees $K = 2/3$ for the direct triple, but the mixed windows require
both the quartic redistribution and the quark-sector angle shift to reach
the observed percent-level agreement.%
\footnote{At each boundary of the interleaving window
($\nu/\mu = q_i^2 q_j^2$, etc.), exactly one sliding window
achieves $K \approx 2/3$; the four boundary values cycle through
the four windows.}%
\footnote{[TO BE REMOVED] At the permutation point
$\nu/\mu = q_1^2 q_2^2$, the unpaired mass is
$d_3 = \mu(q_1 q_2/q_3)^2 = 0.072\MeV$; a seesaw
$(d_3)^2/\Lambda_F = (0.072)^2/102 \approx 0.05\,\mathrm{eV}
\approx m_\nu(\mathrm{atm})$, with $\Lambda_F$ the family
confinement scale.  Single-number coincidence at a non-physical
$\rho$ (ratio $\rho_{\rm phys}/\rho_{\rm perm} \approx 694$).}

\paragraph{Hierarchy of deviations.}
By Eq.~(\ref{eq:correction}), the quartic correction
$\delta M_i/M_i \propto 1/m_i$ is family-dependent:
the heaviest family ($q_3 \approx 1$) is least perturbed.
This is consistent with the waterfall pattern in
Table~\ref{tab:waterfall}: the windows involving the
heaviest quarks---$(t,b,c)$ at $0.37\%$---are
closer to $K = 2/3$ than those involving lighter
quarks---$(b,c,-s)$ at $1.2\%$.

\paragraph{Look-elsewhere effect.}
The waterfall involves four overlapping windows
(consecutive triples from a fixed ordering), each
with a sign choice (up to $2^2 = 4$ independent
patterns per triple after absorbing the overall sign).
A crude trials estimate is
$\binom{6}{3} \times 4 = 80$ for a single window;
since the four windows share quarks and are not
independent, the effective number of independent
trials is smaller than $80^4$.
In the companion letter~\cite{letter_q2} a systematic
Monte Carlo analysis gives $p \approx 1/21\,000$ for
obtaining four windows at or below the observed deviations
from random mass spectra, confirming that the pattern is
unlikely to be accidental.

\paragraph{No new charges.}
The interleaving uses the \emph{same} three family charges
$(q_1, q_2, q_3)$ for both branches.
The doublet structure does not introduce new
charge eigenvalues; it doubles the mass spectrum by
providing two eigenvalues per charge.
The waterfall is a distorted ``triplet~$\times$~2,''
not a many-level hierarchy.

%----------------------------------------------------------------------
\section{Summary}

We have shown:
\begin{enumerate}
\item The direct ($m = \mu q^2$) and inverse ($\hat{m} = \nu/q^2$)
  branches of the family gauge theory have \emph{opposite}
  mass orderings in charge (Section~\ref{sec:branches}).
\item Both branches emerge from the confined spectrum's
  composite doublet $(M_i, \adj(M)_i)$, which has the
  mathematical structure of an $\mathcal{N}{=}2$
  hypermultiplet.  An effective quartic
  $\propto g^2/\mu_\Phi$ couples the branches
  and controls the waterfall deviations
  (Section~\ref{sec:N2}).
\item The six unmixed masses automatically interleave into a
  descending ladder when $\nu/\mu$ lies in a calculable
  window; the SM quark masses satisfy this condition
  (Section~\ref{sec:interleaving}).
\item The sign flip in $K(b,c,-s)$ signals the branch
  mismatch in a mixed window; the specific sign assignment
  constrains the quartic mixing
  (Section~\ref{sec:sign}).
\end{enumerate}
The interleaving is the observable signature of
the composite doublet structure: it produces the up-type
quark sector alongside the down-type sector from a single
$\SU(3)_F$ gauge theory, using the same three family charges.
Quantitative agreement with the full SM quark spectrum
requires either family-dependent mixing or a
quark-sector angle shift, both testable predictions.

\paragraph*{Acknowledgments.}
During the preparation of this work the author used Claude
(Anthropic) for algebraic verification, numerical checks,
and drafting assistance.
The author reviewed and edited all content and takes full
responsibility for the publication.

\begin{thebibliography}{99}

\bibitem{Koide:1983qe}
  Y.~Koide,
  Phys.\ Lett.\ B \textbf{120}, 161 (1983).

\bibitem{Rivero:2005}
  A.~Rivero,
  arXiv:hep-ph/0505220 (2005).

\bibitem{Zenczykowski:2012}
  P.~Zenczykowski,
  Phys.\ Rev.\ D \textbf{86}, 117303 (2012).

\bibitem{letter_q2}
  A.~Rivero, ``Quadratic charge--mass relation from gauged
  family symmetry,'' companion letter (2025).

\bibitem{SeibergWitten:1994a}
  N.~Seiberg and E.~Witten,
  Nucl.\ Phys.\ B \textbf{426}, 19 (1994);
  \textbf{431}, 484 (1994).

\bibitem{PDG:2024}
  S.~Navas \textit{et al.} (Particle Data Group),
  Phys.\ Rev.\ D \textbf{110}, 030001 (2024).

\bibitem{APS:1997}
  P.C.~Argyres, M.R.~Plesser, and A.D.~Shapere,
  Nucl.\ Phys.\ B \textbf{483}, 172 (1997),
  arXiv:hep-th/9705069.

\bibitem{Fayet:1976}
  P.~Fayet,
  Nucl.\ Phys.\ B \textbf{113}, 135 (1976).

\bibitem{Fayet:1977}
  P.~Fayet,
  Phys.\ Lett.\ B \textbf{69}, 489 (1977).

\bibitem{CST:2011}
  C.~Cs\'aki, Y.~Shirman, and J.~Terning,
  Phys.\ Rev.\ D \textbf{84}, 095011 (2011),
  arXiv:1106.3074.

\bibitem{Seiberg:1994}
  N.~Seiberg,
  Nucl.\ Phys.\ B \textbf{435}, 129 (1995),
  arXiv:hep-ph/9411149.

\bibitem{CSS:1996}
  C.~Cs\'aki, M.~Schmaltz, and W.~Skiba,
  Phys.\ Rev.\ Lett.\ \textbf{78}, 799 (1997),
  arXiv:hep-th/9610139.

\bibitem{IS:1995}
  K.~Intriligator and N.~Seiberg,
  Nucl.\ Phys.\ Proc.\ Suppl.\ \textbf{45BC}, 1 (1996),
  arXiv:hep-ph/9509227.

\end{thebibliography}

\end{document}
